\section{VE-SDE 与 VP-SDE 的对比分析}

\subsection{两种 SDE 的核心差异}

在深入学习 Probability Flow ODE 之前，让我们系统对比 VE-SDE 和 VP-SDE 的差异，这有助于理解不同扩散模型的设计取舍。

\begin{table}[H]
\centering
\begin{tabular}{lcc}
\toprule
\textbf{特性} & \textbf{VE-SDE (SMLD)} & \textbf{VP-SDE (DDPM)} \\
\midrule
漂移项 $f(\mathbf{x},t)$ & $0$ & $-\frac{1}{2}\beta(t)\mathbf{x}$ \\
扩散系数 $g(t)$ & $\sqrt{\mathrm{d}[\sigma^2]/\mathrm{d}t}$ & $\sqrt{\beta(t)}$ \\
均值衰减 & 无 ($\alpha_t = 1$) & 有 ($\alpha_t < 1$) \\
方差行为 & 单调递增 (exploding) & 保持恒定 (preserving) \\
终态分布 & $\mathcal{N}(0, \sigma_T^2 \mathbf{I})$ & $\mathcal{N}(0, \mathbf{I})$ \\
SNR 衰减方式 & $1/\sigma^2(t)$ & $\alpha_t^2/(1-\alpha_t^2)$ \\
\bottomrule
\end{tabular}
\caption{VE-SDE 与 VP-SDE 的系统对比}
\end{table}

\begin{knowledgebox}{物理直觉}
\textbf{VE-SDE} 的物理图像是：每个数据点"原地不动"，但周围弥漫的噪声雾越来越浓厚，方差不断增大直到完全淹没数据结构。\textbf{VP-SDE} 的物理图像则是一个 Ornstein-Uhlenbeck 过程：数据点同时被一个弹性力拉向原点（均值衰减），并受到随机扰动（噪声注入），两者平衡使得总方差保持不变。
\end{knowledgebox}

\subsection{实验性能对比}

Song et al. (2021) 的实验表明，在 CIFAR-10 数据集上：
\begin{itemize}
    \item 使用 1000 步 DDPM 采样时，FID（Frechet Inception Distance，越低越好）约为 3--4；
    \item 使用 DDIM 只需约 50 步就能达到类似 FID；
    \item 使用改进的 SDE solver（如基于 Predictor-Corrector 方法），10--15 步即可达到较好质量。
\end{itemize}

\begin{warningbox}{FID 并非唯一指标}
FID 同时衡量生成质量和多样性，但不能完全反映人类对图像质量的感知。在实际应用中，还需要考虑生成速度、模式覆盖率（diversity）、以及下游任务的表现。VE-SDE 和 VP-SDE 在不同设置下各有优劣。
\end{warningbox}

\subsection{本章小结}

VE-SDE 和 VP-SDE 代表了两种不同的加噪策略：前者纯扩散无漂移，后者漂移与扩散平衡。理解它们的差异有助于在实际应用中做出合理的设计选择。


